\begin{abstract}
Checkpoint selection is part of evaluation: the reported model is only as reliable as the rule that chose it. We study training objectives that downweight high-loss examples and show that the objective-aligned validation score can choose unreliable checkpoints under shift. The selected checkpoint can look best on that score, and sometimes on mean accuracy, while doing worse on held-out loss, calibration, and tail- or group-sensitive reliability checks. Under a strong P95 suppressive stress test on Camelyon17 ERM, proxy-only selection worsens held-out loss and ECE on $10/10$ seeds, with a worse fixed-bank teacher-defined hard-example diagnostic as supporting evidence. On end-to-end Camelyon17 finetuning, proxy-only selection raises mean held-out accuracy ($0.683 \rightarrow 0.708$) but worsens held-out loss ($0.902 \rightarrow 1.520$), validation-side ECE ($0.0308 \rightarrow 0.0522$), and the same diagnostic ($7.25 \rightarrow 13.48$). On frozen CivilComments, a validation-loss veto calibrated to the baseline checkpoint changes the run-level verdict: the suppressive trajectories admitted under proxy selection are vetoed on all $10/10$ seeds. The claim is scoped: most evidence compares full reporting workflows, where objective, trajectory, and selector change together; selector-only conclusions come only from explicitly labeled same-trajectory selector and admissibility checks. Same-stack controls show the hazard is not unique to soft clipping: hard bootstrapping and generalized cross-entropy produce the same proxy-improves/reliability-worsens pattern, while label smoothing and focal loss often move the tail in the opposite direction. The effect can be transient or persistent and also appears in narrower frozen-text and winsorized-regression support cases. The practical recommendation is a reporting rule, not a new optimizer: report both proxy-selected and validation-loss-selected checkpoints, and treat harmful disagreement on held-out loss, calibration, or the relevant tail/group metric as an unvalidated reliability claim.
\end{abstract}
